\chapter{Linguagem de domínio para exploração de Regressão Simbólica}

\section{Contexto}
A plataforma proposta nesse projeto tem como um dos objetivos principais permitir a manipulação de modelos de regressão simbólica de maneira similar a outras ferramentas já bem conhecidas no mundo da análise de dados. Uma das ferramentas mais utilizadas na Ciência de Dados é a Structured Query Language (SQL), uma linguagem de domínio~\cite{sommerville2011software} que permite a interação com bancos de dados relacionais.

SQL é uma linguagem de domínio (Domain-Specific Language - DSL), ou seja, uma linguagem que busca resolver um problema específico do universo para qual ela foi pensada. No caso do SQL, o problema que ela resolve é a busca e processamento de dados em bancos de dados relacionais. Outras linguagens de domínio são amplamente utilizadas em computação para simplificar tarefas complexas por meio de abstrações declarativas.

Para facilitar o uso da ferramenta \textit{rEGGression}~\cite{rEGGression} e permitir o uso da plataforma proposta por usuários que já tem conhecimento prático em outras ferramentas voltadas para predição de dados que se utilizam de dialetos SQL, além da plataforma, este trabalho propõe uma linguagem de domínio denominada Inference Query Language (IQL), com sintaxe similar ao SQL.

\section{Inference Query Language}
A Inference Query Language (IQL) foi concebida para permitir que usuários possam consultar soluções dentro de um modelo de regressão simbólica de forma declarativa. Em sistemas de regressão simbólica baseados em grafos de igualdade, como o Eggp~\cite{eggp}, o espaço de soluções pode conter milhares de expressões equivalentes ou sub-expressões que podem ser de interesse para o pesquisador.

Assim como o SQL, o IQL é uma linguagem declarativa. Linguagens declarativas são aquelas em que o programador descreve o que ele deseja obter, sem especificar o fluxo de controle ou os passos procedimentais para alcançar esse resultado~\cite{sommerville2011software}. No contexto do IQL, o usuário pode especificar padrões matemáticos, restrições de complexidade ou métricas de erro, e o sistema é responsável por realizar a busca e filtragem no e-graph.

\subsection{Gramática Formal}
A IQL é composta por uma sintaxe inspirada no padrão ANSI SQL. A gramática formal da linguagem, apresentada na Figura~\ref{fig:iql_grammar}, foi especificada em Extended Backus-Naur Form (EBNF) e define todas as produções aceitas pelo parser da plataforma.

\begin{figure}[htbp]
    \caption{Gramática formal da Inference Query Language (IQL) em notação EBNF, definindo a sintaxe completa para consultas declarativas sobre modelos de regressão simbólica baseados em e-graphs}
    \centering
    \begin{tabular}{rcl}
        \toprule
        \textit{query} & ::= & \texttt{SELECT} [\textit{modifier}] [\texttt{DISTRIBUTION}]
            \textit{column-list} \\
        & & \texttt{FROM} \textit{identifier}
            [\textit{where-clause}] \\
        & & [\textit{pattern-clause}]
            [\textit{order-clause}] \\
        & & [\textit{at-least-clause}]
            [\textit{limit-clause}] \\[4pt]
        \textit{modifier} & ::= & \texttt{TOP} \textit{integer}
            $\mid$ \texttt{PARETO} \\[4pt]
        \textit{column-list} & ::= & \textit{column}
            (\texttt{,} \textit{column})* \\[4pt]
        \textit{column} & ::= & \textit{identifier}
            $\mid$ \textit{predict-call} \\[4pt]
        \textit{predict-call} & ::= & \texttt{PREDICT} \texttt{(} \textit{mapping}
            (\texttt{,} \textit{mapping})* \texttt{)} \\[4pt]
        \textit{mapping} & ::= & \textit{identifier} \texttt{=} \textit{number} \\[4pt]
        \textit{where-clause} & ::= & \texttt{WHERE} \textit{condition}
            (\textit{bool-op} \textit{condition})* \\[4pt]
        \textit{condition} & ::= & \textit{identifier} \textit{comparator}
            \textit{number} \\[4pt]
        \textit{comparator} & ::= & \texttt{=} $\mid$ \texttt{<} $\mid$ \texttt{>}
            $\mid$ \texttt{<=} $\mid$ \texttt{>=} \\[4pt]
        \textit{bool-op} & ::= & \texttt{AND} $\mid$ \texttt{OR} \\[4pt]
        \textit{pattern-clause} & ::= & \texttt{PATTERN IS LIKE} \textit{pattern-expr} \\[4pt]
        \textit{pattern-expr} & ::= & \textit{identifier}
            $\mid$ \textit{pattern-expr} \textit{arith-op} \textit{pattern-expr} \\[4pt]
        \textit{arith-op} & ::= & \texttt{+} $\mid$ \texttt{-}
            $\mid$ \texttt{*} $\mid$ \texttt{/} \\[4pt]
        \textit{order-clause} & ::= & \texttt{ORDER BY} \textit{identifier} \\[4pt]
        \textit{at-least-clause} & ::= & \texttt{AT LEAST} \textit{integer} \\[4pt]
        \textit{limit-clause} & ::= & \texttt{LIMIT} \textit{integer} \\[4pt]
        \textit{identifier} & ::= & \textit{letter} (\textit{letter} $\mid$ \textit{digit}
            $\mid$ \texttt{\_})* \\[4pt]
        \textit{number} & ::= & \textit{digit}+ [\texttt{.} \textit{digit}+] \\[4pt]
        \textit{integer} & ::= & \textit{digit}+ \\
        \bottomrule
    \end{tabular}
    \label{fig:iql_grammar}
\end{figure}

A produção raiz \textit{query} exige obrigatoriamente as cláusulas \texttt{SELECT} e \texttt{FROM}; as demais cláusulas são opcionais. A linguagem opera em três modos principais, controlados pelo modificador que sucede a palavra-chave \texttt{SELECT}:

\begin{itemize}
    \item \textbf{TOP \textit{n}}: Seleciona as $n$ melhores expressões do e-graph de acordo com uma métrica de aptidão (\textit{fitness}). Quando nenhum modificador é especificado, o sistema assume \texttt{TOP 10} como valor padrão;
    \item \textbf{PARETO}: Retorna todas as expressões pertencentes à fronteira de Pareto do modelo, priorizando o equilíbrio entre erro e complexidade estrutural;
    \item \textbf{DISTRIBUTION}: Realiza uma análise de distribuição de padrões estruturais presentes no e-graph, retornando os sub-padrões mais frequentes. Nesse modo, apenas as colunas \texttt{pattern}, \texttt{frequency} e \texttt{fitness} são permitidas na lista de seleção.
\end{itemize}

Os campos que podem ser consultados na cláusula \texttt{SELECT} são determinados pelo executor da linguagem e correspondem às propriedades extraídas das expressões armazenadas no e-graph: \texttt{expression}, \texttt{dl}, \texttt{fitness}, \texttt{latex}, \texttt{numpy}, \texttt{parameters}, \texttt{size}, \texttt{pattern}, \texttt{frequency} e \texttt{prediction}. Além desses campos, a IQL suporta a função \texttt{PREDICT}, que permite ao usuário avaliar expressões com valores específicos para as variáveis independentes, como por exemplo \texttt{PREDICT(x0 = 1.5, x1 = 2.0)}.

A cláusula \texttt{WHERE} aceita condições de filtragem sobre um subconjunto de campos numéricos: \texttt{id}, \texttt{size}, \texttt{parameters} e \texttt{cost}. As condições podem ser compostas por meio dos operadores lógicos \texttt{AND} e \texttt{OR}. A cláusula \texttt{PATTERN IS LIKE} permite a busca por padrões estruturais nas árvores de sintaxe das expressões, utilizando expressões aritméticas como referência para correspondência. A cláusula \texttt{ORDER BY} define a métrica utilizada para ordenar os resultados, enquanto as cláusulas \texttt{AT LEAST} e \texttt{LIMIT} controlam, respectivamente, a frequência mínima dos padrões retornados e o número máximo de resultados no modo \texttt{DISTRIBUTION}.

\subsection{Exemplos de Consultas}
Abaixo são apresentados exemplos de como a IQL pode ser utilizada para extrair conhecimento de modelos de regressão simbólica:

\begin{enumerate}
    \item \textbf{Busca por modelos simples na fronteira de Pareto}:
    \begin{lstlisting}[language=Sql]
    SELECT PARETO expression, size FROM my_model WHERE size <= 5
    \end{lstlisting}

    \item \textbf{Filtragem por padrão estrutural (ex: funções periódicas)}:
    \begin{lstlisting}[language=Sql]
    SELECT TOP 10 expression, fitness, size FROM my_model
    PATTERN IS LIKE sin(v0)
    \end{lstlisting}

    \item \textbf{Busca por interações específicas entre variáveis}:
    \begin{lstlisting}[language=Sql]
    SELECT PARETO expression, size FROM my_model
    PATTERN IS LIKE x1 * x2
    \end{lstlisting}

    \item \textbf{Análise de distribuição de padrões}:
    \begin{lstlisting}[language=Sql]
    SELECT TOP 100 DISTRIBUTION pattern, frequency, fitness
    FROM my_model AT LEAST 5 LIMIT 500
    \end{lstlisting}

    \item \textbf{Predição com valores específicos}:
    \begin{lstlisting}[language=Sql]
    SELECT TOP 5 expression, fitness, PREDICT(x0 = 1.5, x1 = 2.0)
    FROM my_model WHERE size <= 10
    \end{lstlisting}
\end{enumerate}

O uso da IQL permite que o usuário aplique o conhecimento prévio sobre o fenômeno físico diretamente no processo de seleção de modelos, algo que é dificultado em ferramentas de regressão tradicionais que retornam apenas a melhor solução numérica disponível.

\section{Execução}
A implementação da IQL é realizada através de um módulo de consulta que se integra ao back-end da plataforma. O módulo é responsável por transformar a sequência de caracteres enviada pelo usuário em uma árvore de sintaxe abstrata (AST), que é então interpretada para realizar a busca no e-graph do modelo de regressão simbólica. A execução da consulta envolve a aplicação de filtros lógicos, avaliação de métricas e correspondência de padrões estruturais, utilizando as funções de manipulação de e-graphs fornecidas pelo rEGGression~\cite{rEGGression}.

O processamento de tokenização é feito via análise léxica utilizando laços de repetição e condições para identificar os componentes da consulta (palavras-chave, identificadores, operadores, etc.). A construção da AST é realizada por meio de uma técnica chamada Pratt Parsing, a qual associa operadores e precedências para construir a estrutura hierárquica da consulta. O Pratt Parsing permite a criação de \textit{parselets}, classes especializadas em converter uma sequência de tokens em um nó específico da árvore de sintaxe abstrata, para cada tipo de token inicial de uma expressão, facilitando a extensão da linguagem com novos operadores ou funcionalidades no futuro.

A Figura \ref{fig:iql_execution_pipeline} apresenta o diagrama completo do pipeline de execução de uma consulta IQL, desde a entrada textual do usuário até o armazenamento dos resultados no banco de dados.

\begin{figure}[htbp]
    \caption{Pipeline de execução de uma consulta IQL, desde a entrada textual do usuário até o armazenamento dos resultados como InferenceResults no banco de dados}
    \centering
    \begin{tikzpicture}[
        scale=0.72, transform shape,
        node distance=1.8cm and 1.4cm,
        stage/.style={
            rectangle, draw=black, thick, rounded corners=4pt,
            minimum height=1.2cm, minimum width=2.8cm,
            text centered, font=\sffamily\small,
            fill=gray!10
        },
        data/.style={
            rectangle, draw=gray!70, thick, rounded corners=2pt,
            minimum height=0.9cm, minimum width=2.2cm,
            text centered, font=\sffamily\footnotesize,
            fill=white, dashed
        },
        external/.style={
            rectangle, draw=black, thick, rounded corners=4pt,
            minimum height=1.2cm, minimum width=2.8cm,
            text centered, font=\sffamily\small,
            fill=gray!25
        },
        arrow/.style={->, >=stealth, thick, draw=gray!80},
        label/.style={font=\sffamily\tiny, text=gray!70}
    ]
        % Row 1: Input -> Tokenizer -> Token List
        \node[data] (input) {\shortstack{Consulta IQL\\(texto)}};
        \node[stage, right=of input] (tokenizer) {Tokenizer};
        \node[data, right=of tokenizer] (tokens) {\shortstack{Lista de\\Tokens}};

        % Row 1: -> Parser -> AST
        \node[stage, right=of tokens] (parser) {\shortstack{Pratt Parser\\+ Parselets}};
        \node[data, right=of parser] (ast) {AST};

        % Row 2: Executor -> rEGGression -> DataFrame -> Results
        \node[stage, below=1.5cm of ast] (executor) {Executor};
        \node[external, left=of executor] (reggression) {\shortstack{rEGGression\\(e-graph)}};
        \node[data, left=of reggression] (dataframe) {\shortstack{DataFrame\\(Pandas)}};
        \node[stage, left=of dataframe] (mapper) {\shortstack{Mapeamento\\de Colunas}};
        \node[data, below=1.5cm of mapper] (results) {InferenceResults};

        % Arrows Row 1
        \draw[arrow] (input) -- (tokenizer);
        \draw[arrow] (tokenizer) -- node[above, label] {análise léxica} (tokens);
        \draw[arrow] (tokens) -- (parser);
        \draw[arrow] (parser) -- node[above, label] {precedência} (ast);

        % Arrow down
        \draw[arrow] (ast) -- node[right, label] {interpretação} (executor);

        % Arrows Row 2
        \draw[arrow] (executor) -- node[above, label] {top/pareto/dist.} (reggression);
        \draw[arrow] (reggression) -- node[above, label] {busca no e-graph} (dataframe);
        \draw[arrow] (dataframe) -- node[above, label] {renomear/filtrar} (mapper);
        \draw[arrow] (mapper) -- node[right, label] {persistência} (results);

    \end{tikzpicture}
    \label{fig:iql_execution_pipeline}
\end{figure}
